\clearpage
\begin{trapbox}
\textbf{易错点1｜只看表达式，不看定义域}
\begin{center}\begin{tikzpicture}[font=\small,>=stealth,line join=round]
\draw[gray!40,rounded corners=3pt] (0,0) rectangle (14.2,8.8);
\node at (7.1,8.200000000000001) {例1：同一个平方公式，输入范围不同，结论不同};
\begin{scope}[shift={(3.5,2.6)},x=1.05cm,y=1.05cm]\draw[->,gray] (-2.2,0)--(2.4000000000000004,0) node[right]{$x$};\draw[->,gray] (0,-0.3)--(0,4.7) node[left]{$y$};\begin{scope}\clip (-2.2,-0.3) rectangle (2.2,4.5);\draw[blue,thick,domain=-2.1:2.1,samples=120] plot(\x,{\x*\x});\end{scope}\fill[orange!85!black] (1,1) circle(2pt) node[above right]{$(1,1)$};\fill[orange!85!black] (-1,1) circle(2pt) node[above left]{$(-1,1)$};\draw[gray,dotted] (-1,1)--(1,1);\end{scope}\node at (3.5,1.45) {$D=\mathbb R$：是偶函数};\begin{scope}[shift={(10.5,2.6)},x=1.05cm,y=1.05cm]\draw[->,gray] (-2.2,0)--(2.4000000000000004,0) node[right]{$x$};\draw[->,gray] (0,-0.3)--(0,4.7) node[left]{$y$};\begin{scope}\clip (-2.2,-0.3) rectangle (2.2,4.5);\draw[blue,thick,domain=0:2.1,samples=120] plot(\x,{\x*\x});\end{scope}\fill[orange!85!black] (1,1) circle(2pt) node[above right]{$(1,1)$};\draw[red!70!black,thick] (-1-.08,1-.08)--(-1+.08,1+.08) (-1-.08,1+.08)--(-1+.08,1-.08);\node[above left,red!70!black] at (-1,1) {缺失};\end{scope}\node at (10.5,1.45) {$D=[0,+\infty)$：非奇非偶};
\node[text width=12.8cm,align=center] at (7.1,.65) {红叉表示该输入没有函数值，不是函数图像上的点。};
\end{tikzpicture}\end{center}



$f(x)=x^2$若只在$x\geq0$上定义，\textbf{不是偶函数}。右边有$2$，左边却没有$-2$，图像不可能关于$y$轴对称。

\medskip
\textbf{易错点2｜定义域对称，就急着判断是奇还是偶}

$f(x)=x+1$在全体实数上有定义，但$f(-x)=1-x$，既不恒等于$f(x)$，也不恒等于$-f(x)$。\textbf{定义域对称只是第一关。}

\medskip
\end{trapbox}
\clearpage
\begin{trapbox}
\textbf{易错点3｜检查一对数字就下结论}
\begin{center}\begin{tikzpicture}[font=\small,>=stealth,line join=round]
\draw[gray!40,rounded corners=3pt] (0,0) rectangle (14.2,9.4);
\node at (7.1,8.8) {易错：一对点等高，整条图像仍可能不对称};
\begin{scope}[shift={(3.6,3.1)},x=1.2cm,y=0.42cm]\draw[->,gray] (-2.2,0)--(2.4000000000000004,0) node[right]{$x$};\draw[->,gray] (0,-3.5)--(0,11.2) node[left]{$y$};\begin{scope}\clip (-2.2,-3.5) rectangle (2.2,11);\draw[blue,thick,domain=-2.2:2.2,samples=120] plot(\x,{\x*\x*\x+\x*\x-\x});\end{scope}\foreach \x/\y in {-1/1,1/1,-2/-2,2/10}{\fill[orange!85!black] (\x,\y) circle(2pt);}\end{scope}\node[align=left,anchor=west] at (7.8,5.0) {$f(x)=x^2+x(x^2-1)$\\[12pt]$f(1)=f(-1)=1$：这一对符合\\[12pt]但 $f(2)=10,\ f(-2)=-2$\\[12pt]所以不是偶函数\\[12pt]一对符合，不能代替所有点};
\node[text width=12.8cm,align=center] at (7.1,.65) {四个橙点来自原函数；纵轴采用压缩比例，以同时呈现全部样本值。};
\end{tikzpicture}\end{center}



$f(x)=x^2+x(x^2-1)$满足$f(1)=f(-1)=1$，但$f(2)=10$、$f(-2)=-2$，所以它不是偶函数。\textbf{证明成立，要对所有$x$检查；推翻结论，找到一对不符合的就够了。}

\medskip
\end{trapbox}
\clearpage
\begin{trapbox}
\textbf{易错点4｜把$f(0)=0$误认为奇函数的充分证据}
\begin{center}\begin{tikzpicture}[font=\small,>=stealth,line join=round]
\draw[gray!40,rounded corners=3pt] (0,0) rectangle (14.2,9.0);
\node at (7.1,8.4) {例4：输入能成对，不等于输出符合奇偶规律};
\begin{scope}[shift={(3.6,2.9)},x=1.15cm,y=1.0cm]\draw[->,gray] (-2.2,0)--(2.4000000000000004,0) node[right]{$x$};\draw[->,gray] (0,-0.6)--(0,4.7) node[left]{$y$};\begin{scope}\clip (-2.2,-0.6) rectangle (2.2,4.5);\draw[blue,thick,domain=-2.2:2.2,samples=120] plot(\x,{\x*\x+\x});\end{scope}\fill[orange!85!black] (1,2) circle(2pt) node[right]{$(1,2)$};\fill[orange!85!black] (-1,0) circle(2pt) node[below left]{$(-1,0)$};\draw[orange!85!black,dotted] (-1,0)--(-1,2)--(1,2);\end{scope}\node[align=left,anchor=west] at (7.6,4.8) {$f(x)=x^2+x$\\[12pt]$D=\mathbb R$ 对称\\[12pt]$f(1)=2,\ f(-1)=0$\\[12pt]两值既不相等，也不互为相反数\\[12pt]$f(0)=0$ 也不能说明它是奇函数};
\node[text width=12.8cm,align=center] at (7.1,.65) {找到一对违例点，就能否定全称结论。};
\end{tikzpicture}\end{center}



$f(x)=x^2+x$有$f(0)=0$，仍然既非奇也非偶。奇函数\textbf{若在$0$处有定义}，才一定有$f(0)=0$；但反过来不成立。

\medskip
\end{trapbox}
\clearpage
\begin{trapbox}
\textbf{易错点5｜认为奇函数一定经过原点}
\begin{center}\begin{tikzpicture}[font=\small,>=stealth,line join=round]
\draw[gray!40,rounded corners=3pt] (0,0) rectangle (14.2,9.4);
\node at (7.1,8.8) {例3：原点不在定义域，也可以关于原点对称};
\begin{scope}[shift={(3.6,4.9)},x=1.05cm,y=0.6cm]\draw[->,gray] (-2.6,0)--(2.8000000000000003,0) node[right]{$x$};\draw[->,gray] (0,-5.5)--(0,5.7) node[left]{$y$};\begin{scope}\clip (-2.6,-5.5) rectangle (2.6,5.5);\draw[blue,thick,domain=-2.6:-0.185,samples=120] plot(\x,{\x+1/\x});\draw[blue,thick,domain=0.185:2.6,samples=120] plot(\x,{\x+1/\x});\end{scope}\draw[orange!85!black,dashed] (-1,-2)--(1,2);\fill[orange!85!black] (-1,-2) circle(2pt) node[above left]{$(-1,-2)$};\fill[orange!85!black] (1,2) circle(2pt) node[below right]{$(1,2)$};\draw[red!65!black] (0,0) circle(2pt);\node[below right,red!65!black] at (0,0) {原点不在图上};\end{scope}\node[align=left,anchor=west] at (7.9,5.3) {$f(x)=x+1/x$\\[12pt]$D=\mathbb R\setminus\{0\}$ 对称\\[12pt]$f(-x)=-f(x)$\\[12pt]奇函数可以没有 $x=0$};
\node[text width=12.8cm,align=center] at (7.1,.65) {红圈只提示原点缺失；曲线在 $x=0$ 两侧趋向无穷，视窗裁切不是端点。};
\end{tikzpicture}\end{center}



$f(x)=\dfrac1x$的定义域是$x\ne0$，它是奇函数，却没有$(0,0)$这个点。应当说：\textbf{奇函数若有$x=0$这个输入，就必须经过原点。}

\medskip
\end{trapbox}
\clearpage
\begin{trapbox}
\textbf{易错点6｜认为图像关于$x$轴对称的曲线就是函数}


\begin{center}
\begin{tikzpicture}[>=stealth,font=\small]
\draw[gray!40,rounded corners] (0,0) rectangle (14,6.8);
\node at (7,6.3) {关于 $x$ 轴对称：非零输出会冲突，零输出反射后仍是自己};
\begin{scope}[shift={(3.1,3.2)},x=1.3cm,y=1cm]
\draw[->,gray] (-1.5,0)--(2.2,0) node[right]{$x$};
\draw[->,gray] (0,-1.6)--(0,1.8) node[above]{$y$};
\draw[orange!80!black,dashed] (1,-1.5)--(1,1.5);
\fill[blue] (1,1) circle(2pt) node[above left]{$(1,1)$};
\fill[red!70!black] (1,-1) circle(2pt) node[below left]{$(1,-1)$};
\node[align=center] at (.3,-2.3) {若同时保留这两个点\\同一输入 $1$ 就有两个输出};
\end{scope}
\begin{scope}[shift={(10.3,3.2)},x=1.3cm,y=1cm]
\draw[->,gray] (-1.7,0)--(2,0) node[right]{$x$};
\draw[->,gray] (0,-1.6)--(0,1.8) node[above]{$y$};
\draw[blue,very thick] (-1.6,0)--(1.7,0);
\draw[orange!80!black,dashed] (1,-1.4)--(1,1.5);
\fill[orange!80!black] (1,0) circle(2pt) node[above left]{$(1,0)$};
\node at (-.95,1.3){$f(x)\equiv0$};
\node[align=center] at (0,-2.3) {镜像点与原来的点重合\\同一输入仍只有输出 $0$};
\end{scope}
\end{tikzpicture}

\end{center}




$y^2=2px$（$p>0$）关于$x$轴对称，但$x>0$时有两个$y$，不是$y$关于$x$的函数。\textbf{处处为$0$的函数是特别的例外}：它关于$x$轴对称，也不违反一个$x$对应一个$y$的规定。

\medskip
\end{trapbox}
\clearpage
\begin{trapbox}
\textbf{易错点7｜认为零函数在任何定义域上都既奇又偶}
\begin{center}\begin{tikzpicture}[font=\small,>=stealth,line join=round]
\draw[gray!40,rounded corners=3pt] (0,0) rectangle (14.2,6.3);
\node at (7.1,5.7) {易错：同为零函数，也要先看定义域};
\draw[->,gray] (1.0,3.5)--(6.0,3.5) node[right]{$x$};\draw[->,gray] (3.5,2.3)--(3.5,5.3) node[left]{$y$};\draw[blue,very thick] (1.5,3.5)--(5.5,3.5);\fill[orange!85!black] (4.5,3.5) circle(2pt);\fill[orange!85!black] (2.5,3.5) circle(2pt);\node[align=center] at (3.5,1.8) {$D=\mathbb R$：既奇又偶};\draw[->,gray] (8.0,3.5)--(13.0,3.5) node[right]{$x$};\draw[->,gray] (10.5,2.3)--(10.5,5.3) node[left]{$y$};\draw[blue,very thick] (10.5,3.5)--(12.5,3.5);\fill[orange!85!black] (11.5,3.5) circle(2pt);\node[align=center] at (10.5,1.8) {$D=[0,+\infty)$：非奇非偶};
\node[text width=12.8cm,align=center] at (7.1,.65) {蓝线只示意视窗内的图像；左图延伸到所有实数，右图仅保留非负输入。};
\end{tikzpicture}\end{center}



规定$f(x)=0$但只允许$x\geq0$，其图像虽在$x$轴上，定义域却不对称，\textbf{仍然既非奇也非偶}。要同时是奇、偶函数，\textbf{还需要定义域对称}。

\medskip
\textbf{易错点8｜认为偶函数只能有偶次幂，奇函数只能有奇次幂}

$f(x)=|x|$是偶函数，但式子里没有$x^2$一类的偶次幂；$f(x)=\dfrac1x$是奇函数，也不是普通的正整数奇次幂多项式。\textbf{按$f(-x)$验证最可靠，别只靠外形猜。}

\medskip
\end{trapbox}
\clearpage
\begin{trapbox}
\textbf{易错点9｜只计算$f(-x)$，忘了它是否真有意义}
\begin{center}\begin{tikzpicture}[font=\small,>=stealth,line join=round]
\draw[gray!40,rounded corners=3pt] (0,0) rectangle (14.2,8.2);
\node at (7.1,7.6) {易错：能写出代换后的式子，不代表每个输入都有意义};
\begin{scope}[shift={(3.6,3.0)},x=1.1cm,y=1.1cm]\draw[->,gray] (-2.8,0)--(3.0,0) node[right]{$x$};\draw[->,gray] (0,-0.5)--(0,3.7) node[left]{$y$};\begin{scope}\clip (-2.8,-0.5) rectangle (2.8,3.5);\draw[blue,thick,domain=-1:2.8,samples=120] plot(\x,{sqrt(\x+1)});\end{scope}\fill[orange!85!black] (2,{sqrt(3)}) circle(2pt) node[above left]{$(2,\sqrt3)$};\draw[red!70!black,thick] (-2,-.1)--(-2,.1);\node[above,red!70!black] at (-2,0) {$-2$ 不可输入};\end{scope}\node[align=left,anchor=west] at (8.1,4.5) {$f(x)=\sqrt{x+1}$\\[12pt]$D=[-1,+\infty)$\\[12pt]$2\in D$，但 $-2\notin D$\\[12pt]写出 $\sqrt{1-x}$ 前，先检查意义};
\node[text width=12.8cm,align=center] at (7.1,.65) {定义域不对称，已经足以判断非奇非偶，无需继续比较输出。};
\end{tikzpicture}\end{center}



例如$f(x)=\sqrt{x+1}$的定义域为$[-1,+\infty)$，$x=2$能代入，而$-2$不能代入。不能仅凭写出$\sqrt{1-x}$就继续讨论奇偶性；\textbf{先检查输入范围}。
\end{trapbox}
